\section{Background}

\textbf{CVSS.} CVSS \cite{first_cvss_v31} scores intrinsic vulnerability characteristics; we
use a normalized base score as the severity feature and a \texttt{cvss\_only} baseline,
noting its limited exploitation-predictive value \cite{allodi2014comparing}.

\textbf{EPSS.} EPSS \cite{jacobs2021epss} estimates near-term exploitation probability and is
the strongest single public exploitation signal; it is asset-agnostic by design,
which asset-context features attempt to complement. It is our primary baseline.

\textbf{CISA KEV.} The KEV catalog \cite{cisa_kev} lists confirmed-exploited CVEs and, under
BOD 22-01 \cite{cisa_bod2201}, assigns federal remediation deadlines used both as
features and as scheduler deadlines.

\textbf{The vulnerability-host pair.} Severity and exploitation likelihood are
CVE-level, but remediation is performed on a host. We adopt the pair (v, h) as the
unit of decision and explanation: the same CVE on a domain controller and a kiosk
are different decisions, which is what makes per-decision audit evidence
meaningful.

\textbf{Remediation capacity.} Teams remediate within bounded windows under blackout,
CAB cadence, and approval gates; any realistic evaluation must rank \textit{and} schedule
under fixed capacity, since the binding constraint often dominates outcomes.

\textbf{Public-sector audit/compliance context.} Frameworks such as NIST SP 800-40
\cite{nist80040r4}, NIST SP 800-53 \cite{nist80053r5}, and CIS Controls \cite{cis_controls_v8}
emphasize documented, reviewable processes. Producing evidence that *supports
compliance review* is a first-class design goal, distinct from claiming compliance.

\textbf{Synthetic benchmarking.} Real agency telemetry is sensitive and rarely
shareable, impeding reproducibility. A deterministic synthetic fleet from
documented parameter distributions, combined with real public feeds, enables open,
repeatable evaluation; we are explicit that this bounds external validity (Section
11).
